\intro
3D-печать все чаще применяют для изготовления прототипов и мелкосерийных изделий~\cite{лысыч2014области,шимохин2019экономическое}. В отдельных производственных задачах аддитивный подход может уменьшить расход материала и срок изготовления, но его экономическая целесообразность зависит от изделия, выбранной операции, затрат на оборудование и подготовку производства~\cite{шимохин2019экономическое}.. При переходе от единичных прототипов к серийному производству сохраняется задача обеспечения стабильности качества и точности результата.

В рамках данной работы рассматривается метод послойного наплавления (Fused Deposition Modeling, FDM). Изделие формируется последовательным нанесением расплавленного материала по слоям в соответствии с управляющей программой. Для FDM-печати характерна высокая чувствительность к параметрам температуры, скорости печати и перемещений, охлаждения, адгезии первого слоя и к состоянию расходного материала. Отклонение даже одного из параметров может вызвать значительное ухудшение качества изделия. Поэтому контроль процесса должен выполняться непосредственно во время печати.

Практика эксплуатации показывает, что в задачах 3D-печати доля брака может быть существенной даже при наличии отлаженного оборудования. Существенная часть отказов проявляется на ранних этапах печати~\cite{kim2020image} и связана с нарушением устойчивости технологического процесса. При отсутствии своевременного обнаружения такие отклонения приводят к перерасходу материала, простою оборудования, необходимости повторного запуска задания и возможному повреждению принтера.

Практические данные для анализа получены на ферме 3D-принтеров компании Stereotech. При сборе статистики~15~мая~2026~года в истории десяти доступных устройств учтено~4\,640 заданий с известным статусом, из которых~1\,237 отменены. Доля отмен составляет~26,7~\%. Причины отмен в сводке не разделены, поэтому полностью данный показатель нельзя приравнивать к доле изделий с дефектами.

Ключевая практическая сложность контроля связана с тем, что визуальные признаки дефектов возникают локально и стремительно развиваются. При ручном наблюдении оператор не может визуально контролировать каждое задание, особенно при параллельной работе нескольких принтеров. Поэтому в работе акцент сделан на автоматизированном видеоконтроле зоны печати и формализованных сценариях реагирования.

В рамках работы разрабатывается экспериментальный 3D-принтер, оснащённый камерой и алгоритмом раннего обнаружения дефектов и реагированием на них во время печати. Данный алгоритм включает получение кадров из зоны печати, их обработку, выявление и локализацию дефектов по изображению, а также передачу информации в управляющую логику принтера. Под реагированием подразумеваются заранее заданные действия, направленные на ограничение развития дефекта и снижение потерь материала и времени.

Объектом исследования является процесс FDM-печати в условиях малосерийного производства и длительной эксплуатации оборудования. Предметом исследования являются методы и программно-аппаратные средства раннего обнаружения дефектов 3D-печати по данным видеоконтроля, а также способы их интеграции в контур управления принтером для оперативного принятия решений.

Практическая значимость работы заключается в создании первой версии системы, которая переводит контроль качества из ручного режима в режим непрерывного видеонаблюдения с автоматизированным обнаружением отклонений. Такой подход позволяет раньше фиксировать дефекты, сокращать потери материала и времени, а также использовать разработанный принтер как воспроизводимую платформу для последующей инженерной доработки.

Работа актуальна, поскольку несвоевременное обнаружение дефектов 3D-печати приводит к потерям материала и времени~\cite{измайлов2020анализ,petsiuk2022towards}. Если после появления дефекта  выполнение задания продолжается, принтер расходует материал без получения требуемого изделия. Ручной контроль требует регулярного участия оператора, а между проверками дефект может оставаться незамеченным. При длительных заданиях такая зависимость от наблюдения ограничивает автономную работу оборудования.

Визуальный контроль приносит практическую пользу, если результат анализа поступает в контур управления и влияет на ход задания~\cite{oleff2021process}. Поэтому система должна связывать обнаружение дефекта по изображению с реакцией принтера.

Целью данной работы является разработка и тестирование первой версии 3D-принтера с системой раннего обнаружения дефектов в процессе печати и механизмом своевременного реагирования. Для достижения поставленной цели необходимо решить следующие задачи:
\begin{enumerate}
\item Проанализировать дефекты FDM-печати и определить целевой дефект для автоматического обнаружения.
\item Обосновать архитектуру аппаратно-программного комплекса и выбрать компоненты управления, камеру и вычислительный узел обработки изображений.
\item Спроектировать механическую конструкцию экспериментального 3D-принтера и его электрическую подсистему.
\item Подготовить программное окружение принтера, настроить его конфигурацию и получение изображений камеры.
\item Сформировать и разметить набор изображений, разделить данные на обучающую, валидационную и тестовую выборки и подготовить аугментации.
\item Обучить модель обнаружения дефекта типа «спагетти» и оценить качество её работы на тестовой выборке.
\item Реализовать серверную обработку изображений, диагностический обмен и отображение результатов обнаружения.
\item Разработать механизм подтверждения дефекта по последовательности кадров и формирования запроса отмены текущего задания.
\end{enumerate}
